% GENERATED FILE. Edit canonical lessons, then run npm run generate:books.
% canonical-source-hash: fnv1a64:7015a224fc04b460
% canonical-lessons: MR-C107-khare, MR-C107-hirva, MR-C107-pivla, MR-C107-lal, MR-C107-nila

\chapter{True, Green, Yellow, Red, Blue}
\label{ch:mr-true-green-yellow-red-blue}
\hlchaptermodality{107}

\begin{chapteropening}
\textbf{By the end of this chapter:} \emph{I can say true, green, yellow, red and blue.}

Complete the last lesson of chapter 107: i can say true, green, yellow, red and blue.
\end{chapteropening}

\section[khare]{\mr{खरे} (khare) — true}

\label{lesson:MR-C107-khare}



\begin{quote}
Before the new one: say the Marathi for thin, then the Marathi for light.
\end{quote}

\subsection*{You'll want to know: \mr{खरे}}
\textbf{\mr{खरे}} — \emph{khare} — \textquotedblleft{}true\textquotedblright{}.

\begin{grammarlens}[title={What you've built}]
One more word for this chapter.
\end{grammarlens}

\subsection*{Your turn}
\begin{itemize}
\raggedright
  \item \emph{Say it:} \emph{khare}
  \item \emph{Say it:} \emph{khare}, once more
  \item \emph{Say it:} say \emph{khare}
  \item \emph{Recall:} say the Marathi for thin, then the Marathi for light, then say \emph{khare} again
\end{itemize}

\begin{tcolorbox}[breakable,colback=teal!4,colframe=teal!35!black,title={Before you move on}]
What does \mr{खरे} mean? (\textquotedblleft{}True\textquotedblright{}.) Say it once more.
\end{tcolorbox}
\section[hirvā]{\mr{हिरवा} (hirvā) — green}

\label{lesson:MR-C107-hirva}



\begin{quote}
Before the new one: say the Marathi for light, then the Marathi for true.
\end{quote}

\subsection*{You'll want to know: \mr{हिरवा}}
\textbf{\mr{हिरवा}} — \emph{hirvā} — \textquotedblleft{}green\textquotedblright{}.

\begin{grammarlens}[title={What you've built}]
One more word for this chapter.
\end{grammarlens}

\subsection*{Your turn}
\begin{itemize}
\raggedright
  \item \emph{Say it:} \emph{hirvā}
  \item \emph{Say it:} \emph{hirvā}, once more
  \item \emph{Say it:} say \emph{hirvā}
  \item \emph{Recall:} say the Marathi for light, then the Marathi for true, then say \emph{hirvā} again
\end{itemize}

\begin{tcolorbox}[breakable,colback=teal!4,colframe=teal!35!black,title={Before you move on}]
What does \mr{हिरवा} mean? (\textquotedblleft{}Green\textquotedblright{}.) Say it once more.
\end{tcolorbox}
\section[pivḷā]{\mr{पिवळा} (pivḷā) — yellow}

\label{lesson:MR-C107-pivla}



\begin{quote}
Before the new one: say the Marathi for true, then the Marathi for green.
\end{quote}

\subsection*{You'll want to know: \mr{पिवळा}}
\textbf{\mr{पिवळा}} — \emph{pivḷā} — \textquotedblleft{}yellow\textquotedblright{}.

\begin{grammarlens}[title={What you've built}]
One more word for this chapter.
\end{grammarlens}

\subsection*{Your turn}
\begin{itemize}
\raggedright
  \item \emph{Say it:} \emph{pivḷā}
  \item \emph{Say it:} \emph{pivḷā}, once more
  \item \emph{Say it:} say \emph{pivḷā}
  \item \emph{Recall:} say the Marathi for true, then the Marathi for green, then say \emph{pivḷā} again
\end{itemize}

\begin{tcolorbox}[breakable,colback=teal!4,colframe=teal!35!black,title={Before you move on}]
What does \mr{पिवळा} mean? (\textquotedblleft{}Yellow\textquotedblright{}.) Say it once more.
\end{tcolorbox}
\section[lāl]{\mr{लाल} (lāl) — red}

\label{lesson:MR-C107-lal}



\begin{quote}
Before the new one: say the Marathi for green, then the Marathi for yellow.
\end{quote}

\subsection*{You'll want to know: \mr{लाल}}
\textbf{\mr{लाल}} — \emph{lāl} — \textquotedblleft{}red\textquotedblright{}.

\begin{grammarlens}[title={What you've built}]
One more word for this chapter.
\end{grammarlens}

\subsection*{Your turn}
\begin{itemize}
\raggedright
  \item \emph{Say it:} \emph{lāl}
  \item \emph{Say it:} \emph{lāl}, once more
  \item \emph{Say it:} say \emph{lāl}
  \item \emph{Recall:} say the Marathi for green, then the Marathi for yellow, then say \emph{lāl} again
\end{itemize}

\begin{tcolorbox}[breakable,colback=teal!4,colframe=teal!35!black,title={Before you move on}]
What does \mr{लाल} mean? (\textquotedblleft{}Red\textquotedblright{}.) Say it once more.
\end{tcolorbox}
\section[niḷā]{\mr{निळा} (niḷā) — blue}

\label{lesson:MR-C107-nila}



\begin{quote}
Before the new one: say the Marathi for yellow, then the Marathi for red.
\end{quote}

\subsection*{You'll want to know: \mr{निळा}}
\textbf{\mr{निळा}} — \emph{niḷā} — \textquotedblleft{}blue\textquotedblright{}.

\begin{grammarlens}[title={What you've built}]
One more word for this chapter.
\end{grammarlens}

\subsection*{Your turn}
\begin{itemize}
\raggedright
  \item \emph{Say it:} \emph{niḷā}
  \item \emph{Say it:} \emph{niḷā}, once more
  \item \emph{Say it:} say \emph{niḷā}
  \item \emph{Recall:} say the Marathi for yellow, then the Marathi for red, then say \emph{niḷā} again
\end{itemize}

\begin{tcolorbox}[breakable,colback=teal!4,colframe=teal!35!black,title={Before you move on}]
What does \mr{निळा} mean? (\textquotedblleft{}Blue\textquotedblright{}.) Say it once more.
\end{tcolorbox}
